\begin{figure}[H]
\centering
\begin{tikzpicture}[font=\small,>=Stealth,scale=1.25]
\draw[->] (-2.0,0)--(2.25,0) node[right] {$u$};
\draw[->] (0,-1.7)--(0,2.0) node[above] {$v$};
\draw[very thick,blue] (0,0) circle (1.5);
\coordinate (P) at ({1.5*cos(38)},{1.5*sin(38)});
\draw[very thick] (0,0)--(P);
\draw[dashed] (P)--({1.5*cos(38)},0);
\draw[dashed] (P)--(0,{1.5*sin(38)});
\fill (P) circle (1.8pt);
\draw[->] (.55,0) arc[start angle=0,end angle=38,radius=.55];
\node at (.75,.24) {$\theta$};
\node[above right] at (P) {$(\cos\theta,\sin\theta)$};
\node[below] at ({1.5*cos(38)},0) {$\cos\theta$};
\node[left] at (0,{1.5*sin(38)}) {$\sin\theta$};
\node[below left] at (-1.6,-1.3) {$C:\ u^2+v^2=1$};
\end{tikzpicture}
\caption{After the Pythagorean identity is proved, the pair $(\cos\theta,\sin\theta)$ lies on the algebraic locus $C$. Familiar coordinate geometry interprets $C$ as the unit circle.}
\label{fig:unit-circle-trig}
\end{figure}
